\section{Computational methods and provenance}

Finite-time $(a,c)$ rasters are used for discovery, not for declaring invariant
bifurcations.  Candidate periodic states are corrected by shooting equations
with an explicit phase condition.  Stability is computed from the variational
flow and nontrivial Floquet multipliers.  Every continuation correction is
checked against the intended minimal-period orbit identity so that a
double-covered parent cannot be mislabeled as a doubled child.

At high period, full-period shooting becomes poorly conditioned and can lose
the flip direction.  We therefore divide an orbit into multiple shooting
segments, enforce continuity between segment endpoints, and obtain the
monodromy action from the product of the segment variational maps.  Child
identity is tested at a common parameter by phase-aligning the entire orbit and
comparing against the double-covered parent, rather than by comparing one
section point.

Numerical tolerances, frozen predictors, failed acceptance gates, hashes, and
output paths are recorded in experiment receipts in the accompanying
repository.  Event predictions used as prospective tests are frozen before the
subsequent crossing scan and refinement.  This separation prevents an
after-the-fact fit from being reported as prediction.

The return-map topology work uses a stricter protocol that is not yet complete.
It must freeze $\Sigma$, $U$, the coordinate or quotient, and the critical-point
criterion before classifying two versus three branches.  The classification
must then survive declared section and coordinate perturbations.  This is the
acceptance boundary for the topology claims below.
